\chapter{Three-Body Numerics And Ejection Statistics}
\label{app:three-body-numerics}

This appendix records the numerical protocol behind the three-body
demonstrations.  The purpose is pedagogical reproducibility.  A simulation in
this repository should specify its equations, units, random sampling rule,
time-step rule, diagnostic invariants, and statistical estimator.  Without
these ingredients, an ejection probability is only a number; with them, it is a
repeatable experiment whose limitations can be discussed.

\section{Hierarchy Of Models}
\label{sec:three-body-model-hierarchy}

The notes use three related models.  Their variables and invariants are not
identical, so the diagnostic printed by a code must match the model being
integrated.

\begin{center}
\begin{tabular}{p{0.30\textwidth} p{0.34\textwidth} p{0.26\textwidth}}
\hline
model & degrees of freedom & main diagnostics\\
\hline
full Newtonian three-body problem &
three massive bodies in inertial space, reduced by translation when desired &
energy, total angular momentum, center-of-mass drift\\
circular restricted three-body problem &
two primaries on a prescribed circular orbit, massless particle in a rotating
frame &
Jacobi integral, zero-velocity curves, Lagrange point locations\\
heliocentric Sun-Jupiter-asteroid ensemble &
Sun and prescribed Jupiter, many noninteracting massless test particles &
loss classification, binned ejection probability, resonance labels\\
binary-single scattering ensemble &
one bound pair and one incoming body in inertial barycentric coordinates &
binding-energy outcome classification, capture/exchange fractions, cross
sections\\
\hline
\end{tabular}
\end{center}

The first model is autonomous and Hamiltonian.  The circular restricted model
is autonomous only after passing to the rotating frame.  The heliocentric
ensemble used for the asteroid lab is a time-dependent restricted model
because Jupiter is prescribed as a moving source.
The binary-single scattering model is again an autonomous Newtonian three-body
problem, but the observable is different: one samples incoming scattering data
and classifies outgoing binding states.

\section{Units And Equations}
\label{sec:three-body-numerical-units}

The asteroid scripts use astronomical units, years, and solar masses:
\[
        G=4\pi^2,\qquad M_\odot=1 .
\]
For a test particle with heliocentric position \(r_i(t)\in\R^2\), velocity
\(v_i(t)\), and prescribed Jupiter position \(r_J(t)\), the equation is
\[
        \ddot r_i
        =
        -GM_\odot \frac{r_i}{|r_i|^3}
        -GM_J\frac{r_i-r_J}{|r_i-r_J|^3}
        -GM_J\frac{r_J}{|r_J|^3}.
\]
The last term is the indirect acceleration that appears because the origin is
attached to the Sun rather than to an inertial barycenter.  Omitting it changes
the physical model as well as the numerical method, so it is exposed as an
explicit code option.

The CR3BP scripts use nondimensional rotating-frame units.  The total primary
mass is \(1\), the primary separation is \(1\), and the rotating frame has
angular velocity \(1\).  If \(\mu\) is the smaller-primary mass fraction, the
primaries sit at
\[
        r_1=(-\mu,0),\qquad r_2=(1-\mu,0).
\]
The rotating-frame particle state is \((x,y,\dot x,\dot y)\), and
\[
        \ddot x-2\dot y=\Omega_x,\qquad
        \ddot y+2\dot x=\Omega_y,
\]
where
\[
        \Omega(x,y)=
        \frac{x^2+y^2}{2}
        +\frac{1-\mu}{r_1(x,y)}
        +\frac{\mu}{r_2(x,y)} .
\]
The Jacobi integral is
\[
        C=2\Omega(x,y)-\dot x^2-\dot y^2 .
\]
The printed value of \(\max_t |C(t)-C(0)|\) is a numerical error diagnostic.

For binary-single scattering demos, the natural inertial coordinates are
barycentric.  The equations are the full Newtonian equations
\[
  \ddot q_i
  =
  \sum_{j\ne i}Gm_j\frac{q_j-q_i}{|q_j-q_i|^3},
  \qquad i=1,2,3.
\]
The center-of-mass position and velocity should be subtracted from the initial
data.  The diagnostic is the drift in total energy
\[
  H=\sum_i\frac12m_i|\dot q_i|^2
  -
  \sum_{i<j}\frac{Gm_im_j}{|q_i-q_j|}
\]
and, when relevant, the drift in total angular momentum.  At the end of a
scattering integration, pair binding energies
\[
  E_{ij}
  =
  \frac12\mu_{ij}|v_i-v_j|^2-\frac{Gm_im_j}{|q_i-q_j|},
  \qquad
  \mu_{ij}=\frac{m_im_j}{m_i+m_j},
\]
identify which pair, if any, remains bound.  Negative \(E_{ij}\) means that the
pair is instantaneously bound in the two-body sense; an outcome classifier must
also check that the third body is receding and that the result is stable over
the stated stopping time.

\section{Initial Ensembles}
\label{sec:three-body-initial-ensembles}

For the asteroid ejection laboratory, the ensemble is defined by independent
draws
\[
        a_0\sim U(a_{\min},a_{\max}),\qquad
        e_0\sim U(0,e_{\max}),\qquad
        M_0\sim U(0,2\pi),\qquad
        \varpi_0\sim U(0,2\pi).
\]
Here \(a_0\) is the initial semimajor axis, \(e_0\) is eccentricity, \(M_0\)
is mean anomaly, and \(\varpi_0\) is periapse angle.  The eccentric anomaly
\(E\) solves Kepler's equation
\[
        M_0=E-e_0\sin E.
\]
The true anomaly \(f\) is then
\[
        f
        =
        2\arctan\left(
        \frac{\sqrt{1+e_0}\sin(E/2)}
             {\sqrt{1-e_0}\cos(E/2)}
        \right),
\]
and the heliocentric Kepler data are
\[
        p=a_0(1-e_0^2),\qquad
        r=\frac{p}{1+e_0\cos f},
\]
\[
        (x_{\rm orb},y_{\rm orb})=(r\cos f,r\sin f),
        \qquad
        (v_{x,{\rm orb}},v_{y,{\rm orb}})
        =
        \sqrt{\frac{GM_\odot}{p}}(-\sin f,e_0+\cos f).
\]
A rotation by \(\varpi_0\) places the orbit in the reference plane.  The
random seed is part of the experiment.  Two reported probabilities are
comparable only when the sampling law, seed policy, time horizon, and loss
thresholds are stated.

For a binary-single scattering laboratory, a minimal ensemble specifies
\[
  m_1,m_2,m_3,\qquad a_{\rm bin},e_{\rm bin},\qquad
  v_\infty,\qquad b,\qquad \text{binary phase}.
\]
In a planar equal-mass classroom version one may take a circular initial binary
and sample the binary phase uniformly.  If impact parameters are sampled
uniformly in area, use
\[
  b=b_{\max}\sqrt u,\qquad u\sim U(0,1),
\]
not \(b\sim U(0,b_{\max})\).  Uniform sampling in \(b\) and uniform sampling in
area estimate different integrals, so the convention must be printed with any
capture cross section.

The symbol \(v_\infty\) denotes the relative speed of the incoming body and the
binary center of mass at infinite separation.  A code that begins at a finite
launch radius \(R_0\) should not use \(v_\infty\) itself as the launch speed.
Let \(\mu_{\rm out}=m_3(m_1+m_2)/(m_1+m_2+m_3)\), and let \(r_{13}\) and
\(r_{23}\) be the initial distances from the incoming body to the two binary
members.  The finite-start speed \(v_{\rm launch}\) should satisfy
\[
  \frac12\mu_{\rm out}v_{\rm launch}^2
  -Gm_3\left(\frac{m_1}{r_{13}}+\frac{m_2}{r_{23}}\right)
  =
  \frac12\mu_{\rm out}v_\infty^2 .
\]
For \(R_0\gg a_{\rm bin}\), this reduces to the familiar monopole estimate
\[
  v_{\rm launch}
  \simeq
  \left(v_\infty^2+\frac{2G(m_1+m_2+m_3)}{R_0}\right)^{1/2}.
\]
The planar demo uses signed impact parameters for the trajectories but reports
the conventional axisymmetric area cross section.  Thus uniform-in-area
sampling estimates \(\int_0^{b_{\max}}2\pi b\,P(b)\,\dd b\), while uniform
sampling in \(b\) requires the corresponding \(2\pi b\) weight.

\section{Loss Classification}
\label{sec:three-body-loss-classification}

At each diagnostic time the particle is classified as alive or lost.  The
current demo uses three absorbing loss classes:
\[
\begin{array}{lll}
\text{ejected} &:& |r|>R_{\rm ej}\quad \text{or}\quad
E_{\rm Kep}>0 \text{ and } |r|>a_J,\\[2mm]
\text{solar collision} &:& |r|<R_\odot,\\[1mm]
\text{Jupiter close encounter} &:& |r-r_J|<R_J^{\rm close}.
\end{array}
\]
The Kepler energy used in the first line is
\[
        E_{\rm Kep}=\frac{1}{2}|v|^2-\frac{GM_\odot}{|r|}.
\]
This energy is not an invariant of the perturbed problem.  It is used only as
a practical escape indicator.  The ejection radius and close-encounter radius
therefore must be printed with the output.

For binary-single scattering the absorbing classes are different.  A compact
classroom classifier can use the final pair energies \(E_{ij}\):
\[
\begin{array}{lll}
\text{flyby} &:& \text{the original pair remains the only bound pair},\\[1mm]
\text{exchange} &:& \text{a different pair is the most bound outgoing pair},\\[1mm]
\text{capture} &:& \text{the incoming body remains bound by the stated
criterion},\\[1mm]
\text{ionization} &:& E_{12},E_{13},E_{23}\ge0,\\[1mm]
\text{collision/close encounter} &:& |q_i-q_j|<R_{\rm close}
\text{ for some pair}.
\end{array}
\]
The labels depend on the initial state.  If bodies \(1\) and \(2\) are the
initial binary and body \(3\) is incoming, then a final bound pair \((1,3)\) or
\((2,3)\) is naturally called an exchange.  In a capture problem where an
initially unbound object becomes bound to a binary or to one component, the
same final binding-energy calculation may be summarized as a capture event.
The code must state which language it is using.
In conservative Newtonian dynamics such a capture label is a finite-horizon
binding-energy classification.  A temporarily bound triple can later dissolve,
so permanent capture requires either an additional stability criterion, a much
longer integration protocol, dissipation, or interaction with further bodies.

\section{Time Integration And Invariant Drift}
\label{sec:three-body-time-integration-protocol}

For a separable autonomous Hamiltonian \(H(p,q)=T(p)+V(q)\), the velocity
Verlet method
\[
\begin{aligned}
        p_{m+1/2}&=p_m-\frac{\Delta t}{2}\nabla V(q_m),\\
        q_{m+1}&=q_m+\Delta t\,\pd_pT(p_{m+1/2}),\\
        p_{m+1}&=p_{m+1/2}-\frac{\Delta t}{2}\nabla V(q_{m+1})
\end{aligned}
\]
is symplectic and time reversible.  The asteroid demo uses the same
kick-drift-kick idea in a time-dependent restricted field.  Since the
prescribed Jupiter field depends on time, the map is best understood as a
transparent second-order classroom integrator rather than as a final
high-precision celestial mechanics method.

The CR3BP demo uses classical fourth-order Runge-Kutta.  RK4 is useful because
the formula is easy to inspect, but it is not symplectic and does not preserve
the Jacobi integral exactly.  For long-time transport claims, a structure-aware
method and a convergence study are mandatory.

\begin{figure}[htbp]
\centering
\begin{tikzpicture}[
    box/.style={draw, rounded corners=2pt, align=center, minimum width=2.8cm, minimum height=0.85cm},
    arrow/.style={-{Stealth[length=2.5mm]}, thick}
]
\node[box] (model) at (0,0) {equations\\ and units};
\node[box] (sample) at (3.4,0) {initial\\ ensemble};
\node[box] (integrate) at (6.8,0) {time\\ integration};
\node[box] (classify) at (10.2,0) {loss\\ classification};
\node[box] (estimate) at (6.8,-1.9) {probability\\ estimator};
\node[box] (diagnose) at (3.4,-1.9) {invariant\\ diagnostics};
\draw[arrow] (model) -- (sample);
\draw[arrow] (sample) -- (integrate);
\draw[arrow] (integrate) -- (classify);
\draw[arrow] (classify) -- (estimate);
\draw[arrow] (integrate) -- (diagnose);
\draw[arrow] (diagnose) -- (estimate);
\end{tikzpicture}
\caption{A reproducible three-body numerical experiment reports model,
sampling law, integrator, diagnostics, loss rules, and statistical estimator.}
\label{fig:three-body-numerical-pipeline}
\end{figure}

\section{Ejection Probability As A Statistical Estimator}
\label{sec:ejection-probability-estimator}

Let \(X_i\) be the indicator that asteroid \(i\) is ejected by time \(T\):
\[
        X_i=
        \begin{cases}
        1,&\text{particle \(i\) is ejected by time \(T\)},\\
        0,&\text{otherwise}.
        \end{cases}
\]
The Monte Carlo estimator is
\[
        \widehat p_{\rm ej}(T)=\frac{1}{N}\sum_{i=1}^N X_i .
\]
For independent samples from the specified initial ensemble,
\[
        \mathbb{E}\widehat p_{\rm ej}=p_{\rm ej},\qquad
        \operatorname{Var}(\widehat p_{\rm ej})
        =
        \frac{p_{\rm ej}(1-p_{\rm ej})}{N}.
\]
The plug-in standard error is
\[
        \widehat{\operatorname{se}}
        =
        \sqrt{\frac{\widehat p_{\rm ej}(1-\widehat p_{\rm ej})}{N}} .
\]
The same formula applies to a semimajor-axis bin after replacing \(N\) by the
number of particles in that bin.  Bins with few particles have large sampling
error and should not be overinterpreted.

There are three different limits:
\[
        N\to\infty,\qquad
        \Delta t\to0,\qquad
        T\to\infty .
\]
Increasing the number of particles reduces sampling noise.  Decreasing the
time step reduces discretization error.  Increasing the horizon changes the
event being estimated.  These limits answer different questions and should not
be conflated.

If \(T_i\) is the first loss time of particle \(i\), a no-censoring classroom
survival curve is
\[
  \widehat S(t)
  =
  \frac{1}{N}\sum_{i=1}^N{\bf 1}\{T_i>t\}.
\]
This curve is nonincreasing in \(t\).  It is often more informative than a
single terminal probability because it shows whether losses happen in an early
burst or through a long tail.  Kaplan-Meier machinery is unnecessary unless the
experiment introduces censoring beyond the finite integration horizon.

The binary-scattering analogue is an outcome fraction or a cross section.  For
uniform sampling in area on \(0\le b\le b_{\max}\),
\[
  \widehat\sigma_{\rm cap}
  =
  \pi b_{\max}^2\frac{N_{\rm cap}}{N},
  \qquad
  \widehat\sigma_{\rm exch}
  =
  \pi b_{\max}^2\frac{N_{\rm exch}}{N}.
\]
These formulas estimate the cross section for the chosen \(v_\infty\), binary
elements, stopping rule, and classifier.  They should not be quoted without
those inputs.

\section{Bootstrap Intervals And Resonance Windows}
\label{sec:three-body-bootstrap-resonance-windows}

The binomial standard error is appropriate when one reports a single bin with
independently sampled initial conditions.  For plots with many bins, a
bootstrap interval gives a more flexible classroom diagnostic.  Given the
sample
\[
        (a_i,X_i),\qquad i=1,\dots,N,
\]
where \(a_i\) is the initial semimajor axis and \(X_i\) is an ejection
indicator, form bootstrap samples by drawing \(N\) pairs with replacement from
the observed list.  For each resample compute the same binned ejection
fractions.  The empirical \(16\)th and \(84\)th percentiles give a one-sigma
style interval; the \(2.5\)th and \(97.5\)th percentiles give a rough
\(95\) percent interval.  This estimates sampling variability of the chosen
finite-time experiment.  It does not estimate model error, timestep error, or
the uncertainty caused by using a planar restricted model.

For a \(p:q\) interior mean-motion resonance with Jupiter, the nominal
semimajor axis is
\[
        a_{p:q}=a_J\left(\frac{q}{p}\right)^{2/3}.
\]
A resonance-window comparison chooses a width \(\Delta a\) and compares three
sets:
\[
        |a_0-a_{p:q}|<\Delta a,\qquad
        \Delta a\le |a_0-a_{p:q}|<2\Delta a,\qquad
        |a_0-a_{p:q}|\ge 2\Delta a .
\]
For a more dynamical diagnostic one also tracks a resonant angle.  In a planar
restricted problem, a typical resonant angle for an interior asteroid
resonance is
\[
        \phi_{p:q}=p\lambda_J-q\lambda-(p-q)\varpi ,
\]
where \(\lambda_J\) is Jupiter's mean longitude, \(\lambda\) is the asteroid's
mean longitude, and \(\varpi\) is the asteroid periapse longitude.  Libration
of \(\phi_{p:q}\) is stronger evidence of resonant behavior than proximity in
initial semimajor axis alone.  The upgraded demo reports nominal resonance
locations, nearest-resonance labels, finite-window loss fractions, and survival
checkpoints.  A long project version can add full resonant-angle time series;
the quick classroom version keeps the output small and deterministic.

\section{Machine-Readable Output Contracts}
\label{sec:three-body-output-contracts}

The JSON output of the asteroid demo should contain at least:
\[
\begin{array}{ll}
\texttt{model} & \text{equations, units, masses, and integrator},\\
\texttt{integration} & \text{time horizon, timestep, elapsed time},\\
\texttt{ensemble} & \text{sampling law and random seed},\\
\texttt{resonances} & \text{labels, ratios, and nominal locations},\\
\texttt{summary\_by\_bin} & \text{binned loss and ejection fractions},\\
\texttt{summary\_by\_resonance} & \text{finite-window resonance diagnostics},\\
\texttt{survival\_curve} & \text{finite-time survival checkpoints},\\
\texttt{diagnostics} & \text{global counts and numerical metadata}.
\end{array}
\]
Additional keys may appear, but these names are stable enough for tests and lab
notebooks.  The text output can be human-readable; the JSON output is the lab
record.

The JSON output of a binary-scattering demo should contain:
\[
\begin{array}{ll}
\texttt{model} & \text{Newtonian three-body masses and units},\\
\texttt{ensemble} & \text{binary elements, }v_\infty,b_{\max},
  \text{ sampling rule},\\
\texttt{integration} & \text{step size, horizon, and stopping criteria},\\
\texttt{outcome\_counts} & \text{flyby, capture, exchange, ionization,
  collision},\\
\texttt{outcome\_fractions} & \text{the same counts divided by }N,\\
\texttt{cross\_sections} & \text{reported only with the sampling convention},\\
\texttt{diagnostics} & \text{energy drift and close-approach data}.
\end{array}
\]
Tests should verify that fractions lie in \([0,1]\), sum to one over the
exclusive outcome classes, and that cross sections are nonnegative.

\section{Convergence Table Template}
\label{sec:three-body-convergence-template}

A serious numerical report should include a small convergence table.  The
columns should separate the three main effects:
\[
\begin{array}{c|c|c|c|c}
N & \Delta t & T & \widehat p_{\rm ej}(T) &
\text{diagnostic}\\
\hline
N_0 & \Delta t_0 & T_0 & \cdots & \text{baseline}\\
2N_0 & \Delta t_0 & T_0 & \cdots & \text{sampling change}\\
N_0 & \Delta t_0/2 & T_0 & \cdots & \text{timestep change}\\
N_0 & \Delta t_0 & 2T_0 & \cdots & \text{horizon change}
\end{array}
\]
Only the second row tests Monte Carlo variance.  Only the third tests
discretization.  The fourth row changes the physical event from ``ejected by
\(T_0\)'' to ``ejected by \(2T_0\).''  This simple table prevents a common
misreading of the asteroid lab: a longer run is not a more accurate estimate
of the shorter-run probability.

\section{Recommended Reporting Standard}
\label{sec:three-body-reporting-standard}

A publication-quality classroom run should report:
\begin{enumerate}
\item the exact command line or configuration file;
\item the model equations, including whether the indirect term is included;
\item \(N\), \(T\), \(\Delta t\), random seed, and initial sampling law;
\item all loss thresholds;
\item global loss counts and binned counts;
\item binomial standard errors for probabilities;
\item an invariant-drift or convergence diagnostic appropriate to the model;
\item plots with resonance markers, bin widths, and axis units.
\end{enumerate}

The flagship asteroid script now follows this standard closely enough that its
text, CSV, and JSON outputs can be used as lab records.  For a binary-scattering
run, replace the semimajor-axis bins by the impact-parameter sampling rule,
outcome classifier, binding-energy table, and cross-section estimator.  A more
advanced version of either experiment could add symplectic extended-phase
integration, adaptive encounter handling, and comparison against a
high-accuracy reference integrator.
